% Chapter 2, Figure 47: the convention for alpha_0 and alpha_1 in the damping experiment (scan:
% figures/ch2/fig47.png). One rocket, pivoted on the wind axis at its C.G. (origin), drawn in the original
% deflected position (nose down by alpha_0, solid) and at the greatest overshoot (nose up by alpha_1, a
% phantom line); both angles are measured from the wind axis and taken as positive. Lengths are the scan's
% pixels (150 per inch) from the pivot, drawn 1.25 times the printed size (as Ch2 Fig 39): the rocket of
% Fig 45, 366 long, pivoted 221 behind the nose tip; alpha_0 = 17.6 and alpha_1 = 12.7 degrees as printed.
\documentclass[11pt]{standalone}
\usepackage{tamrfig}
\begin{document}
\begin{tikzpicture}[x=0.0212cm, y=0.0212cm]
  \def\azero{17.55}\def\aone{12.74}\def\nose{221}\def\R{261.6}
  % \therocket: nose tip at the origin, axis along +x, in tens of pixels (the ogive's arithmetic overflows
  % TeX in pixel units)
  \newcommand{\therocket}{%
    \rocketoutline{5.4/0.925, 36.6/0.925}
    \rocketfins{36.6}{0.925}{4.6}{1.9}{3.55}{2.7}}
  % the position at greatest overshoot: the outline style set to the house phantom line, in ink
  \begin{scope}[shift={(180-\aone:\nose)}, rotate=-\aone, scale=10, outline/.style={phantom, draw=ink}]
    \therocket
  \end{scope}
  % the original deflected position
  \begin{scope}[shift={(180+\azero:\nose)}, rotate=\azero, scale=10]
    \therocket
  \end{scope}
  % the wind axis and the rockets' axes, each extended ahead of the nose to its angle
  \draw[centerline] (-333,0) -- (292,0) node[anchor=south east, inner sep=0pt, yshift=4pt, text=ink] {Wind axis};
  \foreach \a in {-\aone, \azero} {
    \draw[centerline] (180+\a:\nose) -- (\a:210);
    \draw[thin line] (180+\a:\nose) -- (180+\a:276);
  }
  % the angles, from the wind axis
  \draw[angle arc] (180:\R) arc[start angle=180, end angle=180-\aone, radius=\R]
    node[pos=0.5, fill=white, inner sep=1pt] {$\alpha_1$};
  \draw[angle arc] (180:\R) arc[start angle=180, end angle=180+\azero, radius=\R]
    node[pos=0.5, fill=white, inner sep=1pt] {$\alpha_0$};
  % names of the positions
  \node[anchor=north west, align=left, inner sep=0pt] at (-187,110) {Position at\\ greatest\\ overshoot};
  \node[anchor=north west, align=left, inner sep=0pt] at (-189,-75) {Original\\ deflected\\ position};
\end{tikzpicture}
\end{document}
